\documentclass[11pt]{article}
\usepackage[margin=1in]{geometry}
\usepackage{amsmath,amssymb,amsthm}
\usepackage[T1]{fontenc}
\title{A cardinality obstruction to a height-gap value set\\
  A disproof of one assertion of conjecture 00000007766}
\author{earthking11}
\date{}

\begin{document}
\maketitle

Conjecture 00000007766 asserts, among other things, that the value set of
a real-valued height difference $D$ on periodic points in
$\mathbb P^1(\overline{\mathbb Q})$ consists of $0$ and every real number at
least a certain positive constant $c$. We show that no real-valued function
on those periodic points can have such a value set.

The algebraic closure $\overline{\mathbb Q}$ is countable: each algebraic
number is a root of a nonzero polynomial in $\mathbb Q[X]$, there are
countably many such polynomials, and each has finitely many roots. Also
\[
  \mathbb P^1(\overline{\mathbb Q})
  =\{[z:1]:z\in\overline{\mathbb Q}\}\cup\{[1:0]\},
\]
so this projective line is countable. For any self-map $F$ of it, the set
$\operatorname{Per}(F)$ of periodic points is a subset of a countable set.
Consequently, for any function $D:\mathbb P^1(\overline{\mathbb Q})\to
\mathbb R$, the image $D(\operatorname{Per}(F))$ is countable.

On the other hand, $[c,\infty)$ is uncountable for every real $c$: the map
$x\mapsto c+e^x$ injects $\mathbb R$ into it. Hence
\[
  [c,\infty)\not\subseteq D(\operatorname{Per}(F)),
\]
and in particular
$D(\operatorname{Per}(F))\neq\{0\}\cup[c,\infty)$.
This directly contradicts the conjecture's stated value-set assertion,
regardless of the precise formula for $D$.

\paragraph{Formalization.}
The Lean file takes $K=\operatorname{AlgebraicClosure}(\mathbb Q)$ and
represents $\mathbb P^1(K)$ by \texttt{Option K}, with \texttt{some z}
representing $[z:1]$ and \texttt{none} representing $[1:0]$.
It proves this type countable from the algebraicity of every element of $K$.
For arbitrary $F$ and arbitrary real-valued $D$, the theorem
\texttt{TLMC7766.no\_claimed\_value\_set} rules out the asserted exact
value set. Its axiom audit contains only \texttt{propext},
\texttt{Classical.choice}, and \texttt{Quot.sound}.

\paragraph{Scope.}
The argument does not require the formula for the canonical heights or the
dual map. It refutes the universal coverage of a real interval claimed for
their difference on algebraic periodic points. It makes no claim about the
separate tail asymptotic.

\end{document}
